\begin{tikzpicture}
    \begin{scope}
        \def\L{4}           % near point
        \def\f{1.25}           % focal length
        \def\h{0.6}         % object height
        \def\p{1.8*\f}      % object-lens distance
        % Compute angles
        \def\tanalpha{\h / (\p - \f)}
        \def\tanbeta{\h / (\p)}
        % Draw the converging lens
        \draw[cyan, line join=round,fill=cyan!15] (0.025,-1.25) coordinate (L1) arc (-30:30:1 and 2.5) coordinate (L2) -- ++(-0.05,0) arc (150:210:1 and 2.5) -- cycle;
        % Draw the axial line and focal points
        \draw ({-2.3*\f}, 0) -- ({2.3*\f}, 0);
        \filldraw[black] (-\f, 0) circle (0.035);
        \filldraw[black] (\f, 0) circle (0.035);
        % Draw the object beyond the focal point
        \draw[ultra thick, -Latex] (-\p, 0) coordinate (O1) -- (-\p, \h) coordinate (O2);
        % Draw rays from the object to the image through the lens
       % \draw[red] (O2) -- (0,\h) -- (\f,0) -- ({(\f*\tanalpha)/(\tanbeta)}, {-\f*\tanalpha});
       % \draw[red] (O2) -- (-\f,0) -- (0,{-\f*\tanalpha}) -- ({(\f*\tanalpha)/(\tanbeta)}, {-\f*\tanalpha});
       % \draw[red] (O2) -- (0,0) -- ({(\f*\tanalpha)/(\tanbeta)}, {-\f*\tanalpha});
       % \draw[ultra thick, violet, -Latex] ({(\f*\tanalpha)/(\tanbeta)}, 0) -- ({(\f*\tanalpha)/(\tanbeta)}, {-\f*\tanalpha});
    \end{scope}
    %
    \begin{scope}[shift={(0,-3.5)}]
        \def\L{4}           % near point
        \def\f{1.25}           % focal length
        \def\h{0.6}         % object height
        \def\p{0.8*\f}      % object-lens distance
        % Draw the converging lens
        \draw[cyan, line join=round,fill=cyan!15] (0.025,-1.25) coordinate (L1) arc (-30:30:1 and 2.5) coordinate (L2) -- ++(-0.05,0) arc (150:210:1 and 2.5) -- cycle;
        % Draw the axial line and focal points
        \draw ({-2.3*\f}, 0) -- ({2.3*\f}, 0);
        \filldraw[black] (-\f, 0) circle (0.035);
        \filldraw[black] (\f, 0) circle (0.035);
        % Draw the object before the focal point
        \draw[ultra thick, -Latex] (-\p, 0) coordinate (O1) -- (-\p, \h) coordinate (O2);
    \end{scope}
    %
    \begin{scope}[shift={(8,0)}]
        \def\L{4}           % near point
        \def\f{1.25}           % focal length
        \def\h{0.6}         % object height
        \def\p{1.8*\f}      % object-lens distance
        % Draw the diverging lens
        \draw[cyan, line join=round, fill=cyan!15] (0.2,1.25) coordinate (L1) arc (150:210:1 and 2.5) -- ++(-0.4,0) coordinate (L2) arc (-30:30:1 and 2.5) -- cycle;
        % Draw the axial line and focal points
        \draw ({-2.3*\f}, 0) -- ({2.3*\f}, 0);
        \filldraw[black] (-\f, 0) circle (0.035);
        \filldraw[black] (\f, 0) circle (0.035);
        % Draw the object beyond the focal point
        \draw[ultra thick, -Latex] (-\p, 0) coordinate (O1) -- (-\p, \h) coordinate (O2);
    \end{scope}
    %
    \begin{scope}[shift={(8,-3.5)}]
        \def\L{4}           % near point
        \def\f{1.25}           % focal length
        \def\h{0.6}         % object height
        \def\p{0.4*\f}      % object-lens distance
        % Draw the diverging lens
        \draw[cyan, line join=round, fill=cyan!15] (0.2,1.25) coordinate (L1) arc (150:210:1 and 2.5) -- ++(-0.4,0) coordinate (L2) arc (-30:30:1 and 2.5) -- cycle;
        % Draw the axial line and focal points
        \draw ({-2.3*\f}, 0) -- ({2.3*\f}, 0);
        \filldraw[black] (-\f, 0) circle (0.035);
        \filldraw[black] (\f, 0) circle (0.035);
        % Draw the object beyond the focal point
        \draw[ultra thick, -Latex] (-\p, 0) coordinate (O1) -- (-\p, \h) coordinate (O2);
    \end{scope}
\end{tikzpicture}